% Auto-generated by generate_tex_fragments.py -- do not hand-edit.
\begin{longtable}{@{}L{2.0cm}L{2.0cm}L{2.1cm}L{6.2cm}L{1.9cm}@{}}
\caption{P1: A governance layer improves decision control beyond content-only education. Record, design, population/setting, finding as charted (verbatim from the library record), and direction relative to the proposition. n=44 linked records (kg/graph.json edges).}\label{tab:p1}\\
\toprule
\textbf{Record} & \textbf{Design} & \textbf{Population\slash{}setting} & \textbf{Finding (as charted)} & \textbf{Direction}\\
\midrule
\endfirsthead
\multicolumn{5}{l}{\textit{P1 continued}}\\
\toprule
\textbf{Record} & \textbf{Design} & \textbf{Population\slash{}setting} & \textbf{Finding (as charted)} & \textbf{Direction}\\
\midrule
\endhead
\bottomrule
\endfoot
Abramsky T 2014 \cite{ref24} & Pair-matched cluster randomized controlled trial (8 \seqsplit{communities}) & Community members aged 18-49, Kampala, Uganda; baseline n=1,583, follow-up n=2,532 & SASA! was associated with \seqsplit{significantly} lower social acceptance of IPV among women (aRR 0.54, 95\% CI 0.38-0.79) and 52\% lower past-year physical IPV among women (0.48, 95\% CI 0.16-1.39, crossing 1), plus higher likelihood of women receiving supportive community responses to disclosed violence. & Supports\\
\\[3pt]
Alghamdi M 2022 \cite{srp34253076} & \seqsplit{Qualitative} dominant design with \seqsplit{supplemental} \seqsplit{quantitative} data (n=8) & Canadian Muslim immigrant women, 6 countries of origin & Pre-migration adverse childhood \seqsplit{experiences} and family history of IPV, combined with post-migration language\slash{}social-isolation barriers and \seqsplit{internalized} \seqsplit{patriarchal} values, shaped help-seeking and the eventual decision to leave abusive \seqsplit{relationships} among 8 Muslim immigrant women. & Supports\\
\\[3pt]
Blackwood E 2007 \cite{srp17457732} & Historical\slash{}discourse review & Indonesia, state\slash{}religious\slash{}activist discourse, 1980-2000s & Traces how Indonesian state and Islamic discourse shifted from linking normative gender to \seqsplit{heterosexual} marriage toward more direct legal regulation of sexuality via criminal-code proposals, \seqsplit{illustrating} how marriage has functioned as a governance mechanism in national policy discourse. & Supports\\
\\[3pt]
Boyce SC 2023 \cite{srp37316890} & Item response modelling (Partial Credit Model) on population-based dyadic survey & 559 married adolescent-girl-husband dyads (ages 13-18 wives), rural Niger, 2019 & A 5-item, 2-\seqsplit{dimensional} social-norms scale (\seqsplit{challenging} husband authority; general IPV \seqsplit{acceptability}) showed \seqsplit{reliability} and validity; higher '\seqsplit{challenging} husband authority' scores were \seqsplit{statistically} associated with husband \seqsplit{perpetration} of IPV. & Supports\\
\\[3pt]
Böttcher Bettina 2025 \cite{srp41613074} & Cross-cultural validation study & Austrian student population (including men) & The RAS 3-factor structure was supported; decision-making and freedom-from-coercion subscales had \seqsplit{satisfactory} internal \seqsplit{consistency}, while the \seqsplit{communication} subscale showed marginally lower \seqsplit{reliability}; use in men and gender-diverse \seqsplit{individuals} was noted as \seqsplit{constrained}, requiring adaptation. & Supports\\
\\[3pt]
Couture-Carron A 2020 \cite{srp31896308} & \seqsplit{Exploratory} \seqsplit{qualitative} study & South Asian Muslim \seqsplit{communities} (diaspora) & Cultural norms \seqsplit{constructing} dating as shameful were found to intensify fear of parental\slash{}community reaction and strengthen attachment to abusive partners, with direct \seqsplit{implications} for disclosure, help-seeking, and the ability to end abusive \seqsplit{relationships}. & Supports\\
\\[3pt]
GBD 2023 IPV and SVAC \seqsplit{Collaborators} 2026 \cite{ref1} & GBD modelling study (systematic analysis, \seqsplit{spatiotemporal} Gaussian process regression + burden-of-proof risk \seqsplit{methodology}) & 204 countries and \seqsplit{territories}, females aged 15+ (IPV) and males\slash{}females (sexual violence against children), 1990-2023 & In 2023, an estimated 608 million (95\% UI 518-724) females aged 15+ had ever been exposed to IPV, \seqsplit{contributing} 18.5 million (8.74-30.0) DALYs; IPV ranked as the 4th leading risk factor for DALYs among females aged 15-49. Anxiety disorders and major depressive disorder were the leading IPV-attributed health outcomes (5.43 million and 3.96 million DALYs \seqsplit{respectively}). & Supports\\
\\[3pt]
Gausman J 2022 \cite{srp36550423} & \seqsplit{Qualitative} \seqsplit{comparative} study, 8 focus group \seqsplit{discussions} (4 per country) & Syrian refugees (Jordan) and Rohingya refugees (Bangladesh), fathers\slash{}mothers\slash{}\seqsplit{adolescents} & In both displaced Muslim-majority \seqsplit{populations}, \seqsplit{participants} described child marriage norms reinforced across \seqsplit{generations}, with girls themselves describing having little agency in marriage decisions, alongside contextual resistance-to-change dynamics. & Supports\\
\\[3pt]
Gausman Jewel 2023 \cite{srp37922257} & \seqsplit{Psychometric} validation on cross-sectional multi-country survey data & Women in Argentina (n=890), Ghana (n=1,114), and India (n=1,130) & Adapted the Mothers Autonomy in Decision-Making (MADM) scale to \seqsplit{contraceptive} \seqsplit{counselling} contexts, producing the FP-ADM scale, and validated its \seqsplit{psychometric} properties across three culturally distinct LMIC settings including a South Asian one (India). (Verifier correction recorded in the library entry.) & Supports\\
\\[3pt]
Han X 2024 \cite{srp38970860} & Cross-sectional, random-effects Poisson regression (n=201) & Married Salar (Turkic-Muslim minority) women, rural China & Among 201 married Salar Muslim women, 65.6\% reported past-year child neglect, 30.8\% reported past-year IPV, and 37.6\% had \seqsplit{experienced} child marriage; child marriage and IPV \seqsplit{victimization} were associated with downstream child-neglect risk. & Supports\\
\\[3pt]
Hawkins AJ 2012 \cite{ref33} & Meta-analysis (148 evaluation reports) & Mixed university\slash{}laboratory and religious-setting MRE programs, largely US & Moderate-dosage programs (9-20 contact hours) had stronger effects than low-dosage (1-8 hours); a \seqsplit{communication}-skills emphasis \seqsplit{strengthened} \seqsplit{communication} outcomes but not \seqsplit{satisfaction}; no evidence that formal \seqsplit{institutionalization} (manuals, trainer \seqsplit{certification}) improved effects, and no \seqsplit{significant} difference between religious-setting and university\slash{}laboratory delivery. & Supports\\
\\[3pt]
Hazar Seda 2024 \cite{srp38339841} & Cross-cultural adaptation and \seqsplit{psychometric} validation (online survey, EFA\slash{}CFA) & 308 married women of \seqsplit{reproductive} age (15-49), Türkiye, recruited via WhatsApp\slash{}Instagram & The \seqsplit{Reproductive} Autonomy Scale retained its 3-subscale (14-item) structure in Turkish; subscale Cronbach's alpha ranged from 0.64-0.92 (total scale 0.66), with the decision-making subscale showing the lowest \seqsplit{reliability} among subscales tested. (Verifier correction recorded in the library entry.) & Supports\\
\\[3pt]
Hill AL 2019 \cite{srp31306331} & Secondary cross-sectional analysis of baseline survey data & 550 sexually active female high school students, USA (school health center attendees, ages 14-19) & 12\% reported recent \seqsplit{reproductive} coercion and 17\% reported adolescent \seqsplit{relationship} abuse; both were associated with \seqsplit{differences} in care-seeking and sexual health behaviors, measured via validated RC and \seqsplit{relationship}-abuse \seqsplit{instruments}. (Verifier correction recorded in the library entry.) & Supports\\
\\[3pt]
Islam MK 2021 \cite{srp34527969} & Mixed methods: cross-sectional survey (n=96) + \seqsplit{qualitative} interviews (n=18) + provider interviews (n=9) & Rohingya adolescent girls (age 10-19) with experience of child marriage, Cox's Bazar refugee camp, Bangladesh & Mean age at first marriage was 15.7 years; over 80\% had given birth or were pregnant within two years of the survey; key drivers of child marriage included perceived physical\slash{}mental maturity norms, social norms, insecurity, and family honour. & Supports\\
\\[3pt]
Islam MK 2021 \cite{ref5} & Cross-sectional survey, \seqsplit{multivariable} logistic regression (n=486) & Rohingya refugee women, Cox's Bazar settlement, Bangladesh & 61.32\% of surveyed women had \seqsplit{experienced} child marriage; child marriage was associated with \seqsplit{significantly} higher odds of justifying husband-\seqsplit{perpetrated} physical abuse (AOR=2.71, 95\%CI 1.78-4.11) and of \seqsplit{experiencing} IPV in the past 12 months (AOR=1.72, 95\%CI 1.13-2.61). & Supports\\
\\[3pt]
Jacubowski N 2008 \cite{ref112} & \seqsplit{Documentary}\slash{}archival analysis and expert-informant interviews & Married women, Indonesia (Muslim-majority context) & \seqsplit{Traditional} practices including polygamy, early marriage and contract marriage (mut'a) were identified as enhancing married Indonesian women's likelihood of HIV \seqsplit{acquisition} within marriage, \seqsplit{challenging} the assumption that marital status is protective. & Supports\\
\\[3pt]
Kurnaz Kaan 2026 \cite{srp42015126} & Cross-sectional \seqsplit{descriptive}-\seqsplit{correlational} survey with path analysis & 480 married women aged 18-49, Türkiye & Used the \seqsplit{Reproductive} Coercion Scale, Fertility Desire Scale, and Gender Roles Attitude Scale together; \seqsplit{reproductive} coercion and \seqsplit{traditional} gender-role attitudes were linked to fertility desire pathways in married women (path analysis). (Verifier correction recorded in the library entry.) & Supports\\
\\[3pt]
Kuswanto H 2024 \cite{ref3} & Cross-sectional, binary logistic regression (nationally \seqsplit{representative}, n=9,333) & Females aged 15-20, Indonesia (2017 Indonesia \seqsplit{Demographic} and Health Survey) & National prevalence of female child marriage in Indonesia was 12.53\%; health insurance status and sex of household head were not \seqsplit{significant} predictors, while other \seqsplit{socioeconomic}\slash{}education factors were associated with child marriage risk. & Supports\\
\\[3pt]
McCollum 2008 \cite{ref46} & narrative review of outcome research and clinical practice & IPV-affected couples in conjoint treatment settings, US clinical literature & Concludes conjoint IPV treatment can be effective and safe only when embedded in a broader community response, with careful screening of couples for inclusion, \seqsplit{modification} of standard conjoint methods to promote safety, and ongoing safety assessment with \seqsplit{contingency} plans for escalation; a 'one-size-fits-all' approach is explicitly rejected given IPV typologies. & Supports\\
\\[3pt]
Mieth K 2025 \cite{srp40091116} & \seqsplit{Qualitative} study (49 in-depth interviews, 16 focus group \seqsplit{discussions}) & Forcibly Displaced Myanmar Nationals (Rohingya) aged 15-24, Cox's Bazar, Bangladesh (post-October 2016 arrivals) & \seqsplit{Participants} reported perceived increases in child marriage rates post-\seqsplit{displacement}, attributed to protection fears, \seqsplit{socioeconomic} need, lost education\slash{}employment \seqsplit{opportunity}, and perceived loosening of legal-age \seqsplit{enforcement} in the camp context. & Supports\\
\\[3pt]
Riches Eleanor 2023 \cite{srp36657958} & Cross-sectional online survey validation (classical test theory + CFA) & 826 women of \seqsplit{reproductive} age, United Kingdom & The RAS was highly acceptable (>97.7\% response rate) with good internal \seqsplit{consistency} (Cronbach's alpha 0.75) and fair-good test-retest \seqsplit{reliability} (ICC 0.67); construct validity was supported via \seqsplit{confirmatory} factor analysis. & Supports\\
\\[3pt]
Rogers MM 2022 \cite{srp36693203} & \seqsplit{Qualitative} interview study & Family medicine providers\slash{}staff at two US academic outpatient centers & Notes that one in five men report lifetime IPV \seqsplit{perpetration} and two in three male \seqsplit{perpetrators} seek routine health services, yet family medicine physicians feel unprepared to screen male patients for IPV \seqsplit{perpetration} due to lack of knowledge\slash{}training, despite validated screening tools existing. (Verifier correction recorded in the library entry.) & Supports\\
\\[3pt]
Rumble L 2018 \cite{ref4} & Cross-sectional secondary analysis, \seqsplit{multivariate} regression (nationally \seqsplit{representative}, n=6,578) & Women aged 20-24, Indonesia (2012 Indonesian \seqsplit{Demographic} and Health Survey \slash{} Adolescent \seqsplit{Reproductive} Health Survey) & In a nationally \seqsplit{representative} Indonesian sample, \seqsplit{approximately} 17\% of women were married before age 18 and 6\% before age 16; the paper notes research on child marriage in Southeast Asia \seqsplit{specifically} was scarce prior to this nationally \seqsplit{representative} \seqsplit{multivariate} analysis. & Supports\\
\\[3pt]
Sardinha L 2022 \cite{ref2} & Systematic review + Bayesian multilevel modelling of prevalence data & 366 eligible studies, 2 million women, 161 countries\slash{}areas, ever-partnered women aged 15-49 & Globally, 27\% (UI 23-31\%) of ever-partnered women aged 15-49 have \seqsplit{experienced} physical or sexual, or both, IPV in their lifetime, with 13\% (10-16\%) in the past year; violence starts early, affecting 24\% (UI 21-28\%) of women aged 15-19 and 26\% (23-30\%) aged 20-24. (Verifier correction recorded in the library entry.) & Supports\\
\\[3pt]
Sharifnia A 2025 \cite{srp39425490} & Meta-\seqsplit{ethnography} (33 \seqsplit{qualitative} studies \seqsplit{synthesized}), 6 databases searched & Muslim women, global (multi-country \seqsplit{qualitative} evidence) & Synthesis of 33 \seqsplit{qualitative} studies identified four major themes - experience\slash{}impact of abuse, risk factors, help-seeking, and coping - and found religion functions both as a risk factor for abuse and a barrier to help-seeking, and \seqsplit{simultaneously} as a coping resource for some women. & Supports\\
\\[3pt]
Stern E 2020 \cite{srp31686618} & \seqsplit{Longitudinal} \seqsplit{qualitative} study & \seqsplit{Heterosexual} couples \seqsplit{participating} in the \seqsplit{Indashyikirwa} curriculum, Rwanda & The curriculum, designed to build \seqsplit{relationship} skills and reduce gender-\seqsplit{inequitable} beliefs\slash{}behaviours\slash{}norms \seqsplit{underpinning} IPV, was associated with reported shifts in individual beliefs, couple \seqsplit{relationship} dynamics and gender norms among \seqsplit{participating} couples over time. (Verifier correction recorded in the library entry.) & Supports\\
\\[3pt]
Stern E 2021 \cite{srp32896204} & \seqsplit{Comparative} \seqsplit{qualitative} case study (secondary thematic analysis) & \seqsplit{Participants} across four IPV-prevention programmes, What Works to Prevent Violence against Women and Girls Programme (Ghana, Rwanda, South Africa, Tajikistan) & Similar cross-context pathways of change were identified: learning and applying \seqsplit{relationship} skills, \seqsplit{participatory} challenge of harmful gender norms with group rapport, and \seqsplit{integration} of economic-\seqsplit{empowerment} activities to reduce drivers of IPV and build self-confidence. & Supports\\
\\[3pt]
Taft 2024 \cite{ref92} & randomized controlled trial, n=138 couples & US military service members and partners on a military \seqsplit{installation} & Service members randomized to the couples-based Strength at Home Couples (SAH-C) prevention program showed greater effect-size reductions than a comparison Supportive Prevention condition across physical, sexual, \seqsplit{psychological}, and \seqsplit{specifically} coercive-control IPV, plus greater reductions in \seqsplit{suicidality}; partner-reported effects were in the same direction but less robust. & Supports\\
\\[3pt]
Thombs BD 2017 \cite{srp28789659} & \seqsplit{Comparative} \seqsplit{methodological} review of guideline \seqsplit{organizations} (includes IPV screening as one of 22 compared \seqsplit{recommendations}) & Cross-national comparison of screening guideline bodies (Canada, UK, US) & Among 22 \seqsplit{questionnaire}-based screening \seqsplit{recommendations} compared across three major guideline bodies (including IPV screening), \seqsplit{recommendations} diverged \seqsplit{substantially}, and direct RCT evidence for the \seqsplit{effectiveness} of the \seqsplit{questionnaire}-based screening itself (as opposed to downstream treatment) was often absent or indirect. & Supports\\
\\[3pt]
Upadhyay UD 2014 \cite{ref16} & Instrument \seqsplit{development} and validation (cross-sectional survey, factor analysis) & 1,892 women at 13 family planning and 6 abortion facilities, USA & Generated 26 candidate items, retained 14 items across 3 subscales (freedom from coercion, \seqsplit{communication}, decision-making) via factor analysis; the freedom-from-coercion and \seqsplit{communication} subscales were inversely associated with \seqsplit{unprotected} sex in the past 3 months, supporting construct validity. & Supports\\
\\[3pt]
\seqsplit{Ratnaningsih} M 2022 \cite{srp36163286} & Cross-sectional survey, bivariate\slash{}ANOVA\slash{}logistic regression (n=1,000 households) & Households, Bone District (South Sulawesi) and Palu\slash{}Sigi\slash{}Donggala (Central Sulawesi), Indonesia & Household financial security, dowry practices and local legal frameworks were \seqsplit{significantly} associated with community \seqsplit{acceptability} of child marriage across two Indonesian provinces where child-marriage rates remain above the national average. (Verifier correction recorded in the library entry.) & Supports\\
\\[3pt]
Ghuman SJ 2003 \cite{ref93} & \seqsplit{Comparative} study, 15 matched Muslim\slash{}non-Muslim community pairs & Married women, India, Malaysia, the \seqsplit{Philippines} and Thailand & Women's autonomy was not \seqsplit{consistently} lower in Muslim than non-Muslim \seqsplit{communities} across the four Asian countries studied (including Malaysia and Thailand), and the \seqsplit{association} between autonomy measures and child mortality was weak both across and within \seqsplit{communities}. & Complicates\\
\\[3pt]
Hawkins AJ 2008 \cite{ref6} & Meta-analysis (86 reports, 117 studies, >500 effect sizes) & Mostly US\slash{}Western MRE program \seqsplit{participants}; \seqsplit{predominantly} White, middle-class samples & \seqsplit{Relationship}-quality effect sizes for \seqsplit{experimental} studies ranged d=.30-.36 and \seqsplit{communication}-skills effects d=.43-.45; moderate-dosage programs \seqsplit{outperformed} low-dosage; the authors flag that racial\slash{}ethnic and economic diversity was too thin to draw reliable \seqsplit{conclusions} for \seqsplit{disadvantaged} couples, and that outcomes \seqsplit{policymakers} care about (\seqsplit{relationship} stability, aggression) were rarely measured. & Complicates\\
\\[3pt]
Markman HJ 2013 \cite{ref9} & Randomized controlled trial (couples via religious \seqsplit{organizations} randomized to PREP vs. naturally occurring premarital services; n=44\slash{}66\slash{}83) & US couples marrying through religious \seqsplit{organizations}, 8-year follow-up on divorce & No overall difference in divorce rates between PREP (clergy-delivered or \seqsplit{professional}-delivered) and usual premarital services; couples with higher premarital negative-\seqsplit{communication} observed pre-\seqsplit{intervention} were MORE likely to divorce if given PREP, while low-negativity couples benefited; a premarital history of physical aggression showed a similar (marginally \seqsplit{significant}) harmful-moderation pattern. & Complicates\\
\\[3pt]
Miller Elizabeth 2016 \cite{ref110} & Cluster randomized controlled trial (ARCHES) & 4,009 females aged 16-29 across 25 family planning clinics (17 clusters), USA & The universal education\slash{}counseling \seqsplit{intervention} (ARCHES) did NOT \seqsplit{significantly} reduce \seqsplit{reproductive} coercion (ARR 1.50, 95\% CI 0.95-2.35) or partner violence (ARR 1.07, 95\% CI 0.84-1.38) at 1-year follow-up, though it improved secondary outcomes such as \seqsplit{recognition} and self-efficacy. & Complicates\\
\\[3pt]
Phonkacha N 2026 \cite{srp42300599} & Cross-sectional national survey, binary probit regression & General population, national survey, Thailand (Muslim identity modeled as a covariate) & Muslim identity, poverty, and residence in southern\slash{}central Thailand correlated with higher IPV tolerance for both genders; internet access (as distinct from device ownership) \seqsplit{significantly} reduced IPV tolerance, suggesting a role for \seqsplit{information} access in shifting attitudes. & Complicates\\
\\[3pt]
Bhandari TR 2014 \cite{srp26982668} & Scale \seqsplit{construction} and validation (24-item pool, \seqsplit{psychometric} testing, cross-validation across two districts) & Women in Rupandehi and Kapilvastu districts, Nepal (South Asian, \seqsplit{patriarchal} marital-household context) & Produced a 23-item women's autonomy scale with strong \seqsplit{psychometric} properties (Cronbach's alpha 0.84, test-retest r=0.87, content validity ratio 0.8), validated \seqsplit{specifically} for South Asian household\slash{}marital decision contexts, showing good convergent and \seqsplit{discriminant} validity. (Verifier correction recorded in the library entry.) & Relevant only\\
\\[3pt]
Clyde TL 2020 \cite{ref115} & Narrative review\slash{}position paper & General discussion of \seqsplit{contemporary} premarital-education practice, US-centric & Reviews social trends (\seqsplit{individualism}\slash{}commitment \seqsplit{ambivalence}, changing marriage attitudes, premarital \seqsplit{relationship} histories, media \seqsplit{environment}) and argues current premarital \seqsplit{interventions} need updated content and structure to remain relevant to engaged couples today; offers four proposals for \seqsplit{redesigning} premarital \seqsplit{interventions}. & Relevant only\\
\\[3pt]
Kyegombe N 2017 \cite{srp27682273} & \seqsplit{Qualitative} dyadic study (in-depth interviews, framework analysis) & Couples in SASA! \seqsplit{intervention} \seqsplit{communities}, Kampala, Uganda & Engagement with SASA! activities and \seqsplit{relationship}-skills\slash{}\seqsplit{communication} learning was associated with reduced conflict and IPV for many couples and increased trust, respect and joint economic problem-solving, though \seqsplit{experiences} and degree of change varied widely across couples. (Verifier correction recorded in the library entry.) & Relevant only\\
\\[3pt]
Park J 2026 \cite{srp39757299} & Quasi-\seqsplit{experimental} study with inverse \seqsplit{probability} of treatment weighting (weighted N=291 women, 228 men) & Couples in a city-wide (Seoul) premarital education program, South Korea, during COVID-19 & \seqsplit{Videoconferencing} delivery of the Seoul Premarital Education Program (S-PEP) was associated with increased marital readiness for both men and women, and with increased \seqsplit{relationship} confidence and \seqsplit{satisfaction} for women \seqsplit{specifically}. (Verifier correction recorded in the library entry.) & Relevant only\\
\\[3pt]
Portnoy 2020 \cite{srp32791967} & \seqsplit{qualitative} study --- patient interviews (n=10) and provider focus groups (n=29) & US Veterans Health \seqsplit{Administration} patients and providers & Identifies convergent and divergent patient\slash{}provider beliefs about barriers and \seqsplit{facilitators} to screening \seqsplit{specifically} for IPV \seqsplit{perpetration} (not just \seqsplit{victimization}), an area with a much thinner evidence base than \seqsplit{victimization} screening. & Relevant only\\
\\[3pt]
Semahegn A 2017 \cite{srp29162117} & Study protocol for a community-based quasi-\seqsplit{experimental}\slash{}cluster trial & 1,269 women planned across \seqsplit{intervention}, active comparator and control groups, Awi zone, \seqsplit{northwestern} Ethiopia & The protocol describes a planned pilot of \seqsplit{translating} existing evidence and policy into community-level domestic-violence prevention via peer education, community \seqsplit{representative} training, husband \seqsplit{involvement} and \seqsplit{integration} with community health extension programmes; no outcome data are reported (protocol only). & Relevant only\\
\\[3pt]
US Preventive Services Task Force 2004 \cite{srp14996680} & guideline \slash{} \seqsplit{recommendation} statement & US general population, family and intimate partner violence screening in primary care & USPSTF found the evidence \seqsplit{insufficient} (I statement) to recommend for or against routine screening \seqsplit{instruments} to detect domestic violence, updating a 1996 \seqsplit{insufficient}-evidence finding; explicitly states the letter grade reflects the quality and magnitude of evidence available at that time, not an assumption of no benefit. (Verifier correction recorded in the library entry.) & Relevant only\\
\\[3pt]
Upadhyay UD 2021 \cite{ref18} & Instrument \seqsplit{development} and validation (\seqsplit{qualitative} item generation + national survey, EFA, logistic regression) & 1,117 sexually active young people aged 15-24, USA (national sample) & Developed a 23-item, 7-subscale Sexual and \seqsplit{Reproductive} \seqsplit{Empowerment} Scale for \seqsplit{Adolescents}\slash{}Young Adults (including 'choice of partners, marriage, and children' as a named subscale); validation via logistic regression supported construct validity. & Relevant only\\
\\[3pt]
\end{longtable}
